% Part: normal-modal-logic
% Chapter: frame-correspondence
% Section: properties-accessibility

\documentclass[../../../include/open-logic-section]{subfiles}

\begin{document}

\olfileid{nml}{frd}{acc}

\olsection{د لاسرسي د اړيکو خاصيتونه}

ډېر مودالي !!{formula}s د لاسرسي د اړيکې د ساده او ان اشنا خاصيتونو
ځانګړنې څرګندوي۔ له يوې خوا دا معنا لري چې هر هغه مدل، چې يو ټاکلے
خاصيت ولري، اړوند !!{formula} (او د هغې ټولې تعويضي بېلګې) پکې رښتيا
کوي۔ لومړے د لاسرسي د اړيکو پنځه دوديز ډولونه او هغه فارمولونه ګورو
چې د دغو اړيکو له امله رښتيا وي۔

\begin{thm}\ollabel{thm:soundschemas}
  فرض کړئ $\mModel{M} = \tuple{W, R, V}$ يو مدل دے۔ که $R$ د
  \olref{tab:five} د چپې خوا خاصيت ولري، نو د ښي خوا د !!{formula}
  هره بېلګه په~$\mModel{M}$ کښې رښتيا ده۔
\end{thm}

\begin{table}[t]
    \begin{tabular}{| l || l |}
      \hline
      {\emph{که $R$ \dots وي}} & {\emph{نو \dots په~$\mModel{M}$ کښې رښتيا دے:}} \\
      \hline \hline
      \emph{مسلسل}: $\forall u \exists v Ruv$ & \hbox to.43\textwidth
           {$\Box p \lif \Diamond p$ \hfill (\Ax{D})} \\
      \hline
      \emph{انعکاسي}: $\forall w Rww$  
      & \hbox to.43\textwidth
          {$\Box p \lif p$ \hfill (\Ax{T})} \\
      \hline
      \emph{متناظره}: &  \hbox to .43\textwidth
           {$p \lif \Box\Diamond p$ \hfill (\Ax{B})} \\
           $\forall u\forall v(Ruv \lif Rvu)$ & \\
      \hline
      \emph{انتقالي}: & \hbox to.43\textwidth
           {$\Box p \lif \Box \Box p$ \hfill (\Ax{4})} \\
      $\forall u \forall v \forall w ((Ruv \land Rvw) \lif Ruw)$ & \\
      \hline 
      \emph{اقليدسي}: & \hbox to .43\textwidth
        {$\Diamond p \lif \Box\Diamond p$ \hfill (\Ax{5})} \\
      $\forall w \forall u \forall v ((Rwu \land Rwv) \lif Ruv)$ &\\
      \hline
    \end{tabular}
    \caption{د مطابقت پنځه حقيقتونه۔}
    \ollabel{tab:five}
  \end{table} 


\begin{proof}
  د \Ax{B} مورد داسې دے: په يوه مدل کښې د طرحې د رښتياوالي دپاره
  بايد وښيو چې ټولې بېلګې ئې د مدل په ټولو نړۍو کښې رښتيا دي۔ نو د
  \Ax{B} يوه بېلګه $!A \lif \Box\Diamond !A$ او يوه خپلسرې نړۍ
  $w \in W$ ونيسئ۔ فرض کړئ مقدم~$!A$ په~$w$ کښې رښتيا دے؛ غواړو
  وښيو چې $\Box \Diamond !A$ هم په~$w$ کښې رښتيا دے۔ د دې دپاره بايد
  $\Diamond !A$ له~$w$ څخه د لاسرسي وړ په هرې نړۍ~$w'$ کښې رښتيا
  وي۔ اوس چې د کومې~$w'$ دپاره $Rww'$ وي، د تناظر له فرضه $Rw'w$
  هم راځي (\olref{fig:Bsymm} وګورئ)۔ ځکه چې
  $\mSat{M}{!A}[w]$، نو $\mSat{M}{\Diamond !A}[w']$۔ او ځکه چې
  $w'$ د $Rww'$ لرونکو نړۍو له منځه خپلسره وه، نو
  $\mSat{M}{\Box\Diamond!A}[w]$۔

  نور موردونه د تمرين دپاره پرېږدو۔
\end{proof}

\begin{prob}
  د \olref[nml][frd][acc]{thm:soundschemas} ثبوت بشپړ کړئ۔
\end{prob}

\begin{figure}
  \begin{center}
    \begin{tikzpicture}[modal]
      \node[world] (w1) [label={below:$\mSat{{}}{\formula{A}}$\\
          $\mSat{{}}{\Box\Diamond \formula{A}}$}] {$w$}; 
      \node[world] (w2) [label=below:{$\mSat{{}}{\Diamond \formula{A}}$},
        right=of w1] {$w'$}; 
      \draw[->, bend left] (w1) to (w2); 
      \draw[->, bend left] (w2) to (w1); 
    \end{tikzpicture}
  \end{center}
\caption{له تناظر څخه استدلال۔}
\ollabel{fig:Bsymm}
\end{figure}

پام وکړئ چې د \olref{thm:soundschemas} سرچپه استلزامونه سم نۀ دي:
له دې چې يو مدل يوه طرحه رښتيا کوي، دا نۀ لازمېږي چې د هغه مدل د
لاسرسي اړيکه اړوند خاصيت ولري۔ د \Ax{T} او انعکاسي مدلونو په بحث کښې
داسې مدل په اسانۍ موندل کېږي چې پخپله د \Ax{T} يوه بېلګه پکښې
نارښتيا وي: $W = \{w\}$ او $V(p) = \emptyset$ وټاکئ، او د لاسرسي
اړيکه تشه ونيسئ۔ بيا $R$ انعکاسي نۀ ده، خو
$\mSat{M}{\Box p}[w]$ او $\mSat/{M}{p}[w]$۔ دلته د \Ax{T}
يوازې يوه بېلګه په~$\mModel{M}$ کښې نارښتيا ده؛ نورې بېلګې، لکه
$\Box \lnot p \lif \lnot p$، رښتيا دي۔ د داسې مدل موندل ګران دي
چې د~\Ax{T} \emph{هره تعويضي بېلګه} پکې رښتيا وي او
$\mModel{M}$ انعکاسي نۀ وي۔ خو داسې مدل~$\mModel{M}$ هم شته:
\emph{د سرچينې يادونه: په پورتنۍ يوې نړۍ لرونکې بېلګه کښې اصلي متن
د لاسرسي اړيکه نۀ ټاکي؛ دلته تشه اړيکه څرګنده ټاکل شوې ده۔}

\begin{prop}\ollabel{prop:reflexive}
  فرض کړئ $\mModel{M} = \tuple{W, R, V}$ داسې مدل دے چې
  $W = \{u, v \}$ وي، او نړۍوې $u$ او $v$ د~$R$ له مخې سره تړلې
  وي؛ يعنې $Ruv$ او $Rvu$ دواړه وي، او له دې پرته بله اړيکه نۀ وي۔
  فرض کړئ چې د هر $p$ دپاره $u \in V(p) \Leftrightarrow v
  \in V(p)$۔ بيا:
  \begin{enumerate}
  \item د هر $!A$ دپاره: $\mSat{M}{!A}[u]$ هله او يوازې هله چې
    $\mSat{M}{!A}[v]$ (پر~$!A$ استقرا وکاروئ)۔
  \item د~\Ax{T} هره بېلګه په~$\mModel{M}$ کښې رښتيا ده۔
  \end{enumerate}
  ځکه چې $\mModel{M}$ انعکاسي نۀ دے (بلکې په حقيقت کښې
  \emph{نفي انعکاسي} دے)، د \olref{thm:soundschemas} سرچپه لور
  د \Ax{T} په مورد کښې ناسمه ده (د
  \olref{thm:soundschemas} د ځينو نورو طرحو دپاره هم ورته استدلال
  کېدلے شي، خو د ټولو دپاره نۀ)۔
  \emph{د سرچينې يادونه: اصلي متن يوازې دوه متقابلې څنډې يادوي او
  ځان-څنډې نۀ ردوي؛ د نفي انعکاسيت د دعوې دپاره پورته اړيکه يوازې
  همدغو دوو څنډو ته محدوده شوې ده۔}
\end{prop}

\begin{prob}
  د \olref[nml][frd][acc]{prop:reflexive} دعوې ثابتې کړئ۔
\end{prob}

که څه هم زموږ پام به پر پنځو دوديزو !!{formula}s \Ax{D}،
\Ax{T}، \Ax{B}، \Ax{4} او \Ax{5} وي، په
\olref{tab:anotherfive} کښې د لاسرسي د اړيکو څو نور خاصيتونه هم
راخلو۔ د لاسرسي اړيکه~$R$ هله \emph{جزوي تابعي} ده چې له هرې نړۍ
څخه تر ډېره يوې نړۍ ته لاسرسے وي۔ که له هرې نړۍ څخه کټ مټ يوې نړۍ
ته لاسرسے وي، نو دا اړيکه \emph{تابعي} بولو۔ (نو تابعي اړيکې هماغه
دي چې مسلسل او جزوي تابعي دواړه وي۔) دا اړيکې ځکه «تابعي» بلل
کېږي چې د لاسرسي اړيکه د (جزوي) تابعې په شان کار کوي۔ يوه اړيکه
هله \emph{کمزورې ګڼه} ده چې هر کله $Ruv$ وي، د~$u$ او~$v$ ترمنځ
يوه نړۍ~$w$ وي۔ نو کمزورې ګڼې اړيکې په يوه معنا د انتقالي اړيکو
عکس دي: په انتقالي اړيکه کښې که له~$u$ څخه د~$w$ له لارې په چکر
سره~$v$ ته رسېدے شئ، نو نېغ هم ورته رسېدے شئ؛ په کمزورې ګڼې
اړيکه کښې که له~$u$ څخه نېغ~$v$ ته رسېدے شئ، نو له~$u$ څخه
د کومې~$w$ له لارې په چکر سره هم~$v$ ته رسېدے شئ۔ يوه اړيکه هله
\emph{کمزورې هم‌لوري} ده چې که له يوې نړۍ~$w$ څخه~$u$ او~$v$
دواړو ته رسېدے شئ، نو له~$u$ او~$v$ دواړو څخه يوې ګډې نړۍ~$t$
ته رسېدے شئ؛ دې ته کله کله «الماسي خاصيت» يا «يوځاے کېدنه» هم
وايي۔

\begin{table}[t]
    \begin{tabular}{| p{.50\textwidth} || p{.43\textwidth} |}
      \hline
      {\emph{که $R$ \dots وي}} & {\emph{نو \dots په~$\mModel{M}$ کښې رښتيا دے:}} \\
      \hline \hline
      \emph{جزوي تابعي}: & 
      \multirow{2}{*}{$\Diamond p \lif \Box p$} \\
      $ \forall w \forall u \forall v ((Rwu \land Rwv) \lif u =v )$ & \\
      \hline
      \multirow{2}{*}{\emph{تابعي}: $\forall w \exists v \forall u(Rwu \liff \eq[u][v]) $} 
      & \multirow{2}{*}{$\Diamond p \liff \Box p$} \\
      & \\
      \hline
      \emph{کمزورې ګڼه}: &  \multirow{2}{*}{$\Box \Box p \lif
        \Box p$} \\
      $\forall u\forall v(Ruv \lif \exists w(Ruw \land Rwv))$ & \\
      \hline
      \emph{کمزورې تړلې}: &
      \multirow{3}{*}{\hbox to.43\textwidth{$
        \begin{array}{@{}l@{}}
          \Box ((p \land \Box p) \lif q) \lor {} \\
          \qquad \Box ((q \land \Box q) \lif p)
        \end{array}$ \hfill (\Ax{L})}
      } \\
      $\forall w \forall u \forall v ((Rwu \land Rwv) \lif {}$ &\\
      \qquad $(Ruv \lor u=v \lor Rvu))$ & \\
      \hline 
      \emph{کمزورې هم‌لوري}: & \multirow{3}{*}{\hbox to .43\textwidth
        {$\Diamond\Box p \lif \Box\Diamond p$\hfill (\Ax{G})}} \\
      $\forall w \forall u \forall v ((Rwu \land Rwv) \lif {}$ & \\
      \qquad $\exists t (Rut \land Rvt))$ & \\
      \hline
    \end{tabular}
    \caption{د مطابقت پنځه نور حقيقتونه۔}
    \ollabel{tab:anotherfive}
  \end{table} 

\begin{prob}
  فرض کړئ $\mModel{M} = \tuple{W, R, V}$ يو مدل دے۔ وښيئ چې که
  $R$ د \olref[nml][frd][acc]{tab:anotherfive} د چپې خوا خاصيتونه
  ولري، نو د ښي خوا د هر اړوند فارمول هره بېلګه په~$\mModel{M}$
  کښې رښتيا ده۔
\end{prob}



\end{document}
